% v2-ch12-costs.tex
\chapter{Costs of Access and Continuation}\label{ch:costs}
\spec{How do energy, time, coordination, information, and permission enter as distinct costs?}
{Algebraic path problems (Lehmann); multi-objective shortest paths; weight-constrained paths (Garey and Johnson); list scheduling (Graham).}
{Different costs combine along a path by different operations, and the same graph yields different optimal paths under additive and bottleneck costs (Result~\ref{prop:algebras}); componentwise optima are not jointly achievable (Result~\ref{prop:joint}); scalar weights decide feasibility only for mixtures of paths (Result~\ref{prop:scalarization}); the Pareto front can be exponentially large (Result~\ref{prop:pareto}); integer vector budgets are decided exactly by a product graph (Result~\ref{prop:vector-product}); costs of shared acquisitions are subadditive (Result~\ref{prop:subadd}) and time and energy of a task graph obey different bounds (Result~\ref{prop:time-energy}).}
{None.}
{Chapters~\ref{ch:petri} and~\ref{ch:counting}.}

Chapter~\ref{ch:counting} introduced a budget $\beta$ without saying what it is a budget of. In practice a transition costs several things at once: energy that is spent and may or may not be replenished, time that passes, effort of coordination with others who must also act, information that must be acquired before the transition is possible, and permission that must be obtained. These are not one currency under different names. They combine along a path in different ways, they are consumed from different stocks, and there is no exchange rate fixed by the system that would let one be traded for another. This chapter makes the distinctions precise and records what survives when several costs are carried at once.

\section{Cost algebras}

A \emph{cost algebra} is a semiring $(K,\oplus,\otimes,\mathbf 0,\mathbf 1)$ in which the cost of a path is the product $\otimes$ of its edge costs and the cost of reaching $y$ from $x$ is the sum $\oplus$ over paths. Three examples cover the main cases.
\begin{itemize}[nosep]
  \item \emph{Additive}: $(\mathbb R_{\ge0}\cup\{\infty\},\min,+,\infty,0)$. Energy spent and time elapsed along a single sequence of actions add.
  \item \emph{Bottleneck}: $(\mathbb R_{\ge0}\cup\{\infty\},\min,\max,\infty,0)$. The requirement that is the most demanding along a path governs: the peak power drawn, the highest clearance level required, the slowest of several parties that must agree.
  \item \emph{Boolean}: $(\{0,1\},\vee,\wedge,0,1)$. Permission is available or it is not.
\end{itemize}

\begin{established}[Algebraic path problems]\label{thm:lehmann}
If the semiring is closed (infinite sums of the relevant kind exist and distribute), the path costs between all pairs of an $n$-vertex graph are computed in $O(n^{3})$ semiring operations by the Floyd--Warshall--Kleene recursion. Dijkstra's algorithm computes single-source costs when the semiring is totally ordered and $\otimes$ is monotone.
\end{established}
\source{Lehmann (1977); Tarjan (1981).}

\begin{newresult}[The algebra changes the optimal path]\label{prop:algebras}
There is a graph with two paths from $s$ to $t$ for which the path optimal in the additive algebra is not optimal in the bottleneck algebra and conversely. Under the Boolean algebra the two paths are equally good.
\end{newresult}
\assumes{Edge costs as stated.}
\begin{proof}
Path $P_1$ is a single edge of cost $5$ and path $P_2$ consists of three edges of cost $2$. Under $(\min,+)$ the costs are $5$ and $6$, so $P_1$ is optimal. Under $(\min,\max)$ they are $5$ and $2$, so $P_2$ is optimal. Under the Boolean algebra every edge is permitted and both paths have value $1$.
\end{proof}

Every graph question about cost must therefore state, along with the numbers, how the numbers combine. A requirement stated as a peak, such as a maximum current, is not satisfied by minimizing total energy, and a requirement stated as a total is not satisfied by minimizing a peak.

\section{Several costs at once}

Let each edge carry a cost vector $c(e)\in\mathbb R_{\ge0}^{m}$, with every component additive. A path has cost $c(\pi)=\sum_{e\in\pi}c(e)$, and a vector budget $\beta\in\mathbb R_{\ge0}^{m}$ is met by a path if $c(\pi)\le\beta$ componentwise. The set of vector-budgeted reachable states is $R_\beta(x)=\{y:\exists\pi:x\to y,\ c(\pi)\le\beta\}$. For states $x,y$ let $P(x,y)\subset\mathbb R^m_{\ge0}$ be the set of path costs.

\begin{newresult}[Componentwise optima are not jointly achievable]\label{prop:joint}
Let $\underline c_j(x,y)=\min_{p\in P(x,y)}p_j$ and $\underline c=(\underline c_1,\dots,\underline c_m)$. Then $\underline c\le p$ for every $p\in P(x,y)$, so $\underline c\le\beta$ is necessary for $y\in R_\beta(x)$. It is not sufficient: there is a graph with $P(x,y)=\{(1,5),(5,1)\}$, for which $\underline c=(1,1)\le(2,2)=\beta$ but $y\notin R_\beta(x)$.
\end{newresult}
\assumes{None.}
\begin{proof}
The inequality is the definition of the componentwise minimum. In the example two parallel edges from $x$ to $y$ have costs $(1,5)$ and $(5,1)$, and neither is below $(2,2)$.
\end{proof}

A tempting remedy is to choose weights $w\ge0$ and charge $w\cdot c(e)$ to a single scalar budget. The next result says exactly what that remedy decides.

\begin{newresult}[Scalar weights decide the feasibility of mixtures]\label{prop:scalarization}
Let $P\subset\mathbb R^m_{\ge0}$ be a finite set of path costs and $\beta\in\mathbb R^m_{\ge0}$. Then
\[
\min_{p\in P}w\cdot p\le w\cdot\beta\quad\text{for all }w\ge0
\ \Longleftrightarrow\ 
\beta\in\operatorname{conv}(P)+\mathbb R^m_{\ge0}.
\]
That is, scalar budgets with all weights decide whether some convex combination of path costs meets the vector budget, not whether a single path does. For $P=\{(1,5),(5,1)\}$ the budget $(3,3)$ is the midpoint of the two costs, so every weight vector accepts it, although neither path meets it, whereas the budget $(2,2)$ is rejected by the weight vector $w=(1,1)$.
\end{newresult}
\assumes{Finite $P$.}
\begin{proof}
($\Leftarrow$) If $\beta=\sum_i\lambda_ip_i+z$ with $\lambda$ a probability vector and $z\ge0$ then $w\cdot\beta\ge\sum_i\lambda_iw\cdot p_i\ge\min_pw\cdot p$. ($\Rightarrow$) Suppose $\beta\notin K:=\operatorname{conv}(P)+\mathbb R^m_{\ge0}$. The set $K$ is convex and closed, being the sum of a compact set and a closed cone, so there is $w$ with $w\cdot\beta<\inf_{q\in K}w\cdot q$. If some $w_j<0$ the infimum is $-\infty$ along $q+te_j$, which is impossible, so $w\ge0$. For $w\ge0$ the infimum over $K$ is attained on $P$, so $\min_pw\cdot p>w\cdot\beta$, contradicting the hypothesis.
\end{proof}
Scalarization is thus exactly right when the mission is repeated and costs average, as in the constrained Markov decision processes of Chapter~\ref{ch:mdp}, where randomization over policies realizes the convex combination. It overstates feasibility when the mission is performed once.

\begin{newresult}[The Pareto front can be exponentially large]\label{prop:pareto}
For every $n$ there is a graph with $n+1$ vertices and $2n$ edges and cost vectors in $\mathbb N^{2}$ for which $P(s,t)$ has $2^{n}$ pairwise incomparable elements.
\end{newresult}
\assumes{Two cost components.}
\begin{proof}
Take vertices $v_0,\dots,v_n$ with two parallel edges from $v_i$ to $v_{i+1}$ of costs $(2^{i},0)$ and $(0,2^{i})$, and $s=v_0$, $t=v_n$. A path chooses for each $i$ one of the two edges; if $S$ is the set of indices where the first edge is chosen, the path cost is $(a,2^{n}-1-a)$ with $a=\sum_{i\in S}2^{i}$, and $a$ ranges over $\{0,\dots,2^{n}-1\}$ with each value realized exactly once. All $2^n$ vectors have the same component sum, so if $p\le q$ componentwise then $p=q$, and the vectors are pairwise incomparable.
\end{proof}

The front cannot in general be carried as a state, so a controller that needs to know what it can still afford must carry the remaining budget, as in Chapter~\ref{ch:history-state}.

\begin{newresult}[Exact decision for integer vector budgets]\label{prop:vector-product}
Let costs and the budget be integer vectors. Define the product graph on $X\times\prod_{j=1}^{m}\{0,\dots,\beta_j\}$ with an edge $(x,r)\to(y,r-c(e))$ for each edge $e:x\to y$ with $c(e)\le r$. Then $y\in R_\beta(x)$ if and only if some state $(y,r)$ is reachable from $(x,\beta)$ in the product, and the product has at most $|X|\prod_j(\beta_j+1)$ vertices.
\end{newresult}
\assumes{Integer costs and budgets.}
\begin{proof}
A path in the product from $(x,\beta)$ projects to a path in $X$ whose cost is $\beta-r\le\beta$, where $r\ge0$ is the final remaining budget. Conversely a path in $X$ of cost $p\le\beta$ lifts, step by step, since the remaining budget after each prefix is $\beta$ minus the prefix cost, which is nonnegative.
\end{proof}

The size of the product is polynomial in the numerical values of the budgets and exponential in their length when they are written in binary. This is not an artifact of the construction: deciding whether a path meets two budgets simultaneously is NP-complete, the weight-constrained shortest path problem (Garey and Johnson, ND30), so a vector budget is a source of intractability that a single additive budget is not.

\section{Costs that are shared and costs that overlap in time}

Information and permission are acquired, and once acquired they serve every later use. A cost that is paid for the first use and not again is not additive over a collection of tasks.

\begin{newresult}[Joint acquisition is subadditive]\label{prop:subadd}
In a monotone capability system (Chapter~\ref{ch:endogenous}) with a cost $\kappa(a)\ge0$ for each action, each action being paid for once however often it is used, let $J(T)$ denote the least total cost of a set of actions whose closure from $C_0$ contains the set $T$ of capabilities. Then $J$ is monotone, $J(S\cup T)\le J(S)+J(T)$, and $J(T)\le\sum_{c\in T}J(\{c\})$ with strict inequality possible.
\end{newresult}
\assumes{Monotone capability system with action costs.}
\begin{proof}
If $S\subseteq T$ every set of actions acquiring $T$ acquires $S$, so $J(S)\le J(T)$. The union of an action set acquiring $S$ and one acquiring $T$ acquires both, at cost at most the sum. For strictness take capabilities $c_1,c_2$ each gained by a distinct action of cost $1$ that requires the same capability $c_0$ gained by an action of cost $10$ from $C_0=\varnothing$. Then $J(\{c_1\})=J(\{c_2\})=11$ and $J(\{c_1,c_2\})=12$.
\end{proof}

Computing $J$ is the optimal delete-free planning problem, NP-hard by the result cited in Chapter~\ref{ch:endogenous}. The consequence for accounting is that the cost of a collection of tasks cannot be obtained by summing task by task when tasks share prerequisites, and that the sum over-charges exactly the shared part.

Time differs from energy in how concurrent work is counted. Let a task graph be a directed acyclic graph with durations $d_v>0$ on vertices, let $L$ be the length of a longest path (the critical path), and let $E=\sum_vd_v$ be the total work, which is the energy if each unit of busy time costs one unit of energy.

\begin{newresult}[Time and energy of a task graph]\label{prop:time-energy}
The time to complete the tasks with unlimited parallelism is exactly $L$, while the energy is $E$ however many workers are used. With $p$ identical workers the completion time $T_p$ of any schedule obeys $T_p\ge\max(L,E/p)$, and a greedy list schedule achieves $T_p\le E/p+(1-1/p)L$.
\end{newresult}
\assumes{Tasks are non-preemptive and a task starts only after its predecessors finish.}
\begin{proof}
Each path must be executed in sequence, so no schedule finishes before $L$, and the schedule that starts every task as soon as its predecessors finish takes exactly $L$. $p$ workers perform at most $pT_p$ units of work in time $T_p$, so $T_p\ge E/p$. The upper bound is Graham's analysis of list scheduling (1966).
\end{proof}

The two costs are measured on one graph and optimized by different quantities: time by the longest path, energy by the sum. Adding workers lowers time toward $\max(L,E/p)$ and does not lower energy at all, so a design question of the form ``how much faster for how much more'' cannot be answered by a single exchange rate.


\section{The duality gap of a single mission}\label{sec:gap-mission}

Result~\ref{prop:scalarization} says what scalar weights decide. The same fact can be stated as an optimization bound with a numerical gap. Let a mission have three available paths with energy and time costs $(0,10)$, $(10,0)$ and $(6,6)$, and let the time be limited to $B=6$. The least energy among paths that meet the limit is $6$, attained by the third path. The Lagrangian relaxation prices time at $\lambda\ge0$ and charges $c_1+\lambda(c_2-B)$, giving the lower bound $\max_{\lambda\ge0}\min_p\{p_1+\lambda(p_2-6)\}$. The three terms are $4\lambda$, $10-6\lambda$ and $6$, the inner minimum is the smaller of the first two and $6$, the first two cross at $\lambda=1$ with value $4$, and so the dual value is $4<6$. The gap of $2$ is the price of insisting on a single path: the mixture that takes the second path with probability $0.4$ and the first with probability $0.6$ has expected cost $(4,6)$, meets the limit on average, and attains the dual value $4$.

\begin{newresult}[Lagrangian bound equals the mixture optimum]\label{prop:lagrange}
For a finite set $P\subset\mathbb R^2_{\ge0}$ of path costs and a limit $B$ on the second component, the Lagrangian dual value equals the least expected first component over probability mixtures of paths with expected second component at most $B$, whenever some path has second component at most $B$. The least first component over single paths is at least this value and can exceed it.
\end{newresult}
\assumes{Finite path set, feasibility of at least one single path.}
\begin{proof}
The mixture problem is the linear program $\min\{\sum_i\lambda_ip_{i1}:\sum_i\lambda_i=1,\lambda\ge0,\sum_i\lambda_ip_{i2}\le B\}$. It is feasible by assumption, and a linear program with a feasible bounded problem has zero duality gap; its dual with multiplier $\lambda\ge0$ on the time constraint is exactly $\max_\lambda\min_i\{p_{i1}+\lambda(p_{i2}-B)\}$. Single paths are mixtures with a point mass, so the single-path optimum is at least the mixture optimum. The example above gives $6>4$.
\end{proof}

\section{Adding workers can slow a schedule}\label{sec:anomaly}

Result~\ref{prop:time-energy} bounds the completion time of a greedy schedule but does not claim that it falls as workers are added. For list scheduling it can rise. Take nine tasks of durations $3,2,2,2,4,4,4,4,9$, with $T_1$ preceding $T_9$ and $T_4$ preceding each of $T_5,T_6,T_7,T_8$, and let the list order be $T_1,\dots,T_9$. A simulation of the list rule gives completion time $12$ with three workers and $15$ with four. The critical path is $T_4,T_5$ of length $6$ against $T_1,T_9$ of length $12$, so $L=12$ and the three-worker schedule is optimal, while the four-worker schedule is $25\%$ longer. With three workers, $T_1,T_2,T_3$ start at time $0$ and the worker freed by $T_1$ at time $3$ takes $T_9$ at once, since $T_4,T_5,\dots$ are served earlier but are not yet all ready; $T_9$ then finishes at $12$. With four workers, $T_1,\dots,T_4$ all start at time $0$, and when $T_2,T_3,T_4$ finish at time $2$ the three free workers take $T_5,T_6,T_7$, which precede $T_9$ in the list. The worker freed by $T_1$ at time $3$ takes $T_8$, so $T_9$ cannot start until time $6$ and finishes at $15$. The total work is $34$ either way. The anomaly is that the priority list, which was acceptable for three workers, uses the extra capacity on short tasks and delays the critical task. Extra capacity helps only if the priority rule is critical-path aware, and the example is a concrete reason that the bound $T_p\le E/p+(1-1/p)L$ is an upper bound on the damage and not a description of it.

\section*{Exercises}
\addcontentsline{toc}{section}{Exercises}

\begin{exercise}
Give a graph on which three different paths are optimal in the additive, bottleneck, and fewest-edges senses respectively.
\end{exercise}

\begin{exercise}
For $P=\{(1,5),(5,1),(4,4)\}$ compute the set $\operatorname{conv}(P)+\mathbb R^2_{\ge0}$ and list the budgets $\beta$ accepted by all scalar weights but met by no single path.
\end{exercise}

\begin{exercise}
In the graph of Result~\ref{prop:pareto} with $n=3$ list all eight cost vectors, and give a vector budget accepted by exactly three of them.
\end{exercise}

\begin{exercise}
Show that the product of Result~\ref{prop:vector-product} can be pruned to the states $(y,r)$ that are not dominated by another reachable $(y,r')$ with $r'\ge r$, and explain why the pruned set is the Pareto front of Result~\ref{prop:pareto} in disguise.
\end{exercise}

\begin{exercise}
For the task graph with tasks $a,b,c,d$ of durations $2,3,1,4$ where $d$ depends on $a$ and $b$, and $c$ is independent, compute $L$, $E$, and the bounds of Result~\ref{prop:time-energy} for $p=2$.
\end{exercise}

\begin{exercise}
For the three paths of Section~\ref{sec:gap-mission} compute the dual value for the limit $B=5$ and for $B=7$, and compare it with the best single-path value.
\end{exercise}

\begin{exercise}
Prove that the Lagrangian dual value in Result~\ref{prop:lagrange} is nonincreasing in $B$ and equals the minimum first component over all paths once $B$ is at least the largest second component.
\end{exercise}

\begin{exercise}
Verify by drawing the schedules that the nine-task example of Section~\ref{sec:anomaly} completes at time $12$ with three workers and $15$ with four under the list order $T_1,\dots,T_9$, and find a priority order for which four workers finish at time $12$, and show that no schedule can finish before $12$.
\end{exercise}

\begin{exercise}
Show that a cost algebra in which $\otimes$ is $\max$ and $\oplus$ is $\min$ makes the Pareto front of Result~\ref{prop:pareto} a single point when only the bottleneck on one component is tracked, and explain what this says about which costs create combinatorial difficulty.
\end{exercise}

